\section{The oriented logarithmic action and the exact cubic exponent}
\label{sec:cubic-oriented-action}

We retain the field $F_1$ and the actual cyclotomic automorphism
$\gamma$ of AC-K22.  All matrices below act on columns.
A row defining a character of a quotient therefore transforms
by right multiplication by the action matrix.

\begin{theorem}[AC-K23: the exceptional residue is two]
\label{AC-K23}
On the certified group
$U=\widetilde{\operatorname{Cl}}_3(F_1)\simeq\mathbb Z/9$,
the automorphism fixing $\alpha$ and sending
$\zeta_9+\zeta_9^{-1}$ to $\zeta_9^4+\zeta_9^{-4}$
acts by multiplication by four.
Consequently the exact source exponents are
\[
       \boxed{\displaystyle m_e=4+1_{e\equiv2\pmod3}}
                     \qquad(e\geq1),                 \tag{K23.1}
\]
and the even $K$-theory lengths of AC-K12 are
\[
       \kappa_e=5+2v_3(e)+1_{e\equiv2\pmod3}.         \tag{K23.2}
\]
For $N=t-1\geq0$, the complete fixed-three contribution is
\[
       \sum_{e=1}^{N}m_e\log3
          =\left(4N+\left\lfloor\frac{N+1}{3}\right\rfloor\right)\log3.
                                                        \tag{K23.3}
\]
These assertions are unconditional computer-assisted arithmetic
consequences of the certified class/unit data and the preceding
proved reductions.  They do not use theorem A.
\end{theorem}
\begin{proof}
\proofstep{1. A complete presentation with named divisor classes.}
Put $F_1=\mathbb Q[y]/(p)$, with $p$ as in (KD.5).
Its two primes above three are denoted by $\mathfrak p_1,
\mathfrak p_2$, ordered by local degrees six and three.
Their logarithmic indices, returned by \texttt{bnflogef}, are
$[2,3]$ and $[1,3]$ respectively.
Thus their common logarithmic degree, in the PARI
normalization, is
\[
                          d=3\log_3(4).               \tag{K23.4}
\]
This uses the normalization of Remark~3.6 of
Belabas--Jaulent, \emph{The logarithmic class group package
in PARI/GP}, p.~9 of the author's version, and the documented
\texttt{bnflogdegree} convention.
Their ordinary local degrees are different, so each prime
is fixed by the order-three automorphism $\gamma$.

Let $\mathfrak g_1,\mathfrak g_2$ be the ordinary class
representatives in the certified \texttt{bnf} structure,
with respective orders $21,3$.
Choose the exact elements $t_i$ returned by
\texttt{idealcoprime(b,g[i],3)} and set
$\mathfrak a_i=(t_i)\mathfrak g_i$.
These integral ideals are coprime to three and represent
the same class basis.  Their exact ideal matrices and the
$t_i$ are saved in the certificate; their norms are
$32849$ and $232082623$.
The script verifies that each valuation at each
$\mathfrak p_j$ is zero.
Define
\[
 r_i=\frac{\log_3 N\mathfrak a_i}{d}\in\mathbb Z_3,
 \qquad
 v_j(z)=-\frac{\log_3N_{(F_1)_{\mathfrak p_j}/\mathbb Q_3}(z)}d.
                                                        \tag{K23.5}
\]
The integrality of $r_i$ can also be seen without numerical
approximation.  The field contains the cyclic cubic field
$\mathbb Q(\zeta_9+\zeta_9^{-1})$; the residue degree of
any prime away from three makes its residue norm congruent
to $1$ or $-1$ modulo nine.  Thus the logarithm of every
away-from-three ideal norm is divisible by nine,
as is $d$, with $v_3(d)=2$.
Away from three, the logarithmic divisor coefficients
are the ordinary ideal valuations.

The following exact ideal identities are computed and
verified by equality of ideal Hermite matrices:
\[
 \mathfrak a_i^{d_i}=(u_i),\quad d_1=21,\ d_2=3,
 \qquad
 \mathfrak p_j=(z_j)\mathfrak a_1^{c_{1j}}
                              \mathfrak a_2^{c_{2j}}.
                                                        \tag{K23.6}
\]
Here the exponents and principal generators come from
\texttt{bnfisprincipal}, followed by multiplication by
$\prod_i t_i^{-c_{ij}}$ to express them relative to the
coprime ideals $\mathfrak a_i$ rather than the original
$\mathfrak g_i$.
Let $\varepsilon_1,\ldots,\varepsilon_5$ be the certified
fundamental units.  The unit rank is five and the only
roots of unity in this odd-degree field are $\pm1$;
the latter have zero logarithmic divisor.

In the free $\mathbb Z_3$-module with ordered basis
$(\mathfrak a_1,\mathfrak a_2,\mathfrak p_1,\mathfrak p_2)$,
form the nine columns
\[
 \begin{array}{ll}
 (d_i e_i,\ v_1(u_i),v_2(u_i))^{\mathsf t}, &i=1,2,\\
 (-c_{1j},-c_{2j},\ v_1(z_j),v_2(z_j))^{\mathsf t}, &j=1,2,\\
 (0,0,\ v_1(\varepsilon_k),v_2(\varepsilon_k))^{\mathsf t},
                                                   &1\leq k\leq5.
 \end{array}                                           \tag{K23.7}
\]
In the first row, $d_i e_i$ means the two ordinary-class
coordinates, so each displayed column has four entries.
These columns are principal logarithmic divisors.
They give a complete presentation of the full logarithmic
class group on this named set of generators.
Here is a proof of completeness, needed to compute an action
rather than just abstract invariant factors.

Every prime ideal away from three is an ordinary class
combination of the $\mathfrak a_i$, up to a principal ideal.
Its logarithmic class therefore is a combination of
$\mathfrak a_1,\mathfrak a_2,\mathfrak p_1,\mathfrak p_2$.
Conversely, suppose a combination on this set is the
logarithmic divisor of an element of
$\mathbb Z_3\otimes F_1^\times$.
Its ordinary divisor away from three has the same
$\mathfrak a_i$ coefficients, say $a=(a_1,a_2)^{\mathsf t}$;
write its ordinary valuations at three as $t=(t_1,t_2)^{\mathsf t}$.
Since the ordinary class presentation is exactly
$\operatorname{diag}(21,3)$, there is $k\in\mathbb Z_3^2$ with
\[
                 a+Ct=\operatorname{diag}(21,3)k.
\]
The product $u_1^{k_1}u_2^{k_2}z_1^{t_1}z_2^{t_2}$
then has exactly the same ordinary divisor.
Their quotient belongs to $\mathbb Z_3\otimes\mathcal O_{F_1}^{\times}$,
so it is a combination of the five certified units,
up to torsion with zero logarithmic divisor.
Taking logarithmic divisors expresses the relation as a
combination of (K23.7).
The use of three-adic exponents is justified by tensoring
the ordinary divisor sequence with the flat module
$\mathbb Z_3$.  This exact sequence is printed in
Belabas--Jaulent, \S1, p.~2; the logarithmic definitions
and product formula are on pp.~2--3.
The argument also shows directly why a mere subgroup of
units would not suffice for this action certificate.

\proofstep{2. Pass to degree zero and compute the quotient.}
All columns in (K23.7) have degree zero by the product formula.
A basis of the degree-zero submodule is
\[
 A_1=\mathfrak a_1-r_1\mathfrak p_2,
 \qquad A_2=\mathfrak a_2-r_2\mathfrak p_2,
 \qquad C=\mathfrak p_1-\mathfrak p_2.                 \tag{K23.8}
\]
Indeed a vector with four coordinates has degree divided
by $d$ equal to $r_1x_1+r_2x_2+x_3+x_4$.
In the basis (K23.8), its degree-zero coordinates are
simply its first three coordinates.
Let $M_0$ be the resulting three-by-nine relation matrix.
The certificate gives
\[
 (r_1,r_2)\equiv(7,20)\pmod{27},\qquad
 M_0\equiv
 \begin{pmatrix}
 21&0&26&8&0&0&0&0&0\\
 0&3&0&0&0&0&0&0&0\\
 6&4&24&21&0&0&18&6&6
 \end{pmatrix}\pmod{27}.                             \tag{K23.9}
\]
There is no degree-zero surjectivity assumption on an
$S$-class quotient: (K23.8) is the actual kernel of the
degree map in this complete presentation.
The primitive row
\[
                         q_0=(0,1,6)                 \tag{K23.10}
\]
annihilates $M_0$ modulo nine, as direct multiplication
of (K23.9) verifies.
It therefore defines a surjection from the degree-zero
logarithmic class group to $\mathbb Z/9$.
That group has order nine by AC-K22, so the surjection
is an isomorphism.
As an additional finite check, the Smith invariants of
$(M_0\bmod27\mid27I_3)$ are $[9,1,1]$.

\proofstep{3. Compute the action, preserving its orientation.}
The exact automorphism from AC-K22 fixes $\alpha$ and
sends a root $\beta$ of $X^3-3X+1$ to $2-\beta-\beta^2$.
Since
$\beta=\zeta_9+\zeta_9^{-1}$ implies
$2-\beta-\beta^2=\zeta_9^4+\zeta_9^{-4}$,
this is $\gamma$, not its inverse.
For each $i$, compute and verify the exact ideal identity
\[
       \gamma(\mathfrak a_i)=(\eta_i)
                       \prod_{k=1}^2\mathfrak a_k^{b_{ki}}.
                                                        \tag{K23.11}
\]
In the logarithmic class group this means
\[
       [\gamma(\mathfrak a_i)]
        =\sum_k b_{ki}[\mathfrak a_k]
                    -\sum_j v_j(\eta_i)[\mathfrak p_j].
                                                        \tag{K23.12}
\]
The minus sign follows by setting the principal
logarithmic divisor of $\eta_i$ equal to zero.
Both $\mathfrak p_j$ are fixed.
If $G$ is this four-by-four matrix, its degree-zero
matrix $G_0$ is obtained by applying $G$ to the three
columns (K23.8) and retaining the first three rows.
The result is
\[
                 G_0\equiv
                    \begin{pmatrix}1&14&0\\0&1&0\\0&8&1\end{pmatrix}
                            \pmod{27}.                \tag{K23.13}
\]
In particular
\[
                         q_0G_0=4q_0\pmod9.           \tag{K23.14}
\]
Through the isomorphism (K23.10), the actual action of
$\gamma$ is therefore multiplication by four.
This conclusion is unchanged by a different generator
of the cyclic quotient: its automorphism group is abelian.

\proofstep{4. Certification and arithmetic consequence.}
The exact class and unit data were fully certified in AC-K22.
All ideal equalities in (K23.6) and (K23.11) are verified
by exact arithmetic, rather than by floating-point embeddings.
For each of the two local factors of $p$ over $\mathbb Q_3$,
the script evaluates the norm by a resultant and then
its Iwasawa logarithm.  The factors are uniquely matched
to the primes by their distinct degrees six and three.
The source justification for this finite-precision local
norm calculation is Belabas--Jaulent, Lemma~2.1, p.~6;
its logarithm loses the valuation of the norm, and the
subsequent division by $d$ loses $v_3(d)=2$ further digits.
The actual arithmetic retains absolute precision at least
$76$ in $M_0$ and $77$ in $G_0$ after all these losses.
Thus reduction modulo $27$ or nine in (K23.9)--(K23.14)
is certified with a margin of more than seventy digits.
The computed degree sums vanish to the retained precision;
their exact vanishing follows separately from the product formula.

The complete reproducible script is
\path{evidence/cubic_logarithmic/real_cubic_logarithmic_gamma.gp}.
Its output records all equality checks, both matrices,
their minimum absolute precisions, the quotient row,
its relation test, its action test, the finite Smith check,
and \texttt{CERTIFICATE\_PASS=1}.
The accompanying binary retains the exact class/unit data,
the exact field automorphism, coprime representatives,
principal generators and the local approximations used.
No degree-fifty-four class calculation is used.

Finally AC-K21 and AC-K22 give $r=-1\equiv2\pmod3$,
which proves (K23.1).
The exact identity
$\kappa_e=m_e-1+2(1+v_3(e))$ from AC-K12 gives
(K23.2).
There are exactly $\lfloor(N+1)/3\rfloor$ positive
integers at most $N$ congruent to two modulo three;
this proves (K23.3).
\end{proof}

\paragraph{Interpretation and remaining scope.}
The cyclic-nine logarithmic invariant, with its labelled
cyclotomic action, is enough to determine every fixed-cubic
source exponent.  Its Tate specialization has a longer
quotient exactly for the residue class in (K23.1).
Through AC-K10--AC-K12 these are wild even $K$-theory
localization kernels, not the full even $K$-groups.
In the finite-\'etale coefficient dictionary of
AC-K20, the oriented finite action is additional arithmetic
data on the coefficient module; the categorical dictionary
does not calculate it.  The estimate for all other primes,
the growing-host main Selmer depth and the geometric
identification of the relevant theta character remain separate.
